%% kapa: exam=1 round=26 year=2015 period=2 form=A subject=accounting subject_ko=회계학 seq=17 number=57
%% source: Q-Net 감정평가사 기출문제 (cst00309) artlSeq=5138001 file=extracted/1008/5138001/제26회_감정평가사 1차_기출/제26회_감정평가사 1차_기출/2교시_감평_A형.pdf
%% conversion: kordoc 4.24.51 -> tools/gen1.py (auto); review=proofread-sonnet; flags=none
\begin{question}{57}
(주)감평은 20×1년 1월 1일에 다음과 같은 전환사채를 액면발행하였다.
\begin{framedtext}
○ 액면금액 : ￦500,000 \par
○ 표시이자율 : 연 8\% \par
○ 일반사채의 시장수익률 : 연 10\% \par
○ 이자지급방법 : 매 연도 말 후급 \par
○ 상환기일(만기) : 20×3년 12월 31일 \par
○ 원금상환방법 : 상환할증금은 없으며, 상환기일에 액면금액을 일시에 상환함
\end{framedtext}
현가계수표는 다음과 같다.

{\small\centering\begin{tabular}{|c|c|c|c|c|}\hline
\multirow{2}{*}{구분} & \multicolumn{2}{c|}{\makecell{단일금액\\￦1의 현재가치}} & \multicolumn{2}{c|}{\makecell{정상연금\\￦1의 현재가치}} \\ \cline{2-5}
 & 8\% & 10\% & 8\% & 10\% \\ \hline
1년 & 0.9259 & 0.9091 & 0.9259 & 0.9091 \\ \hline
2년 & 0.8573 & 0.8264 & 1.7833 & 1.7355 \\ \hline
3년 & 0.7938 & 0.7513 & 2.5771 & 2.4868 \\ \hline
\end{tabular}\par}
\medskip
20×1년 1월 1일 전환사채 발행시 부채요소와 자본요소로 계상될 금액은 각각 얼마인가?

\noindent\hspace*{2em}\makebox[8em][l]{\uline{부채요소}}\uline{자본요소}\par
\begin{choices}
  \item \makebox[8em][l]{￦500,000}￦0
  \item \makebox[8em][l]{￦24,878}￦475,122
  \item \makebox[8em][l]{￦4,976}￦495,024
  \item \makebox[8em][l]{￦475,122}￦24,878
  \item \makebox[8em][l]{￦495,747}￦4,253
\end{choices}
\end{question}
\answer{4}
